\chapter{Memory Access Instructions}
\label{ch:memory-instructions}

\section{Overview}

Memory access instructions transfer data between registers and memory. Following the load/store architecture, these are
the only instructions that can access memory—all arithmetic and logical operations must work on registers.

The SIRC-1 provides two primary memory operations:
\begin{itemize}
	\item \mnemonic{LOAD} -- Read from memory into a register
	\item \mnemonic{STOR} -- Write from a register to memory
\end{itemize}

\mnemonic{LOAD} and \mnemonic{STOR} each support four memory addressing forms. Both support indirect addressing with
an immediate offset and indirect addressing with a register offset; \mnemonic{LOAD} additionally supports
post-increment and \mnemonic{STOR} additionally supports pre-decrement:
\begin{itemize}
	\item Indirect with immediate offset: \texttt{(\#offset, addr)}
	\item Indirect with register offset: \texttt{(rO, addr)}
	\item Post-increment (\mnemonic{LOAD} only): \texttt{(\#offset, addr)+} or \texttt{(rO, addr)+}
	\item Pre-decrement (\mnemonic{STOR} only): \texttt{-(\#offset, addr)} or \texttt{-(rO, addr)}
\end{itemize}

\textbf{Important:} \mnemonic{STOR}'s register-offset forms support an optional shift suffix, applied to the source
register value before it is written to memory. \mnemonic{LOAD} does not support a shift suffix on any form: the
offset register \texttt{rO} is always used unshifted. See Section~\ref{sec:memory-shifts} for details. Shifts are
also \textbf{not} supported when loading immediate values or copying register values directly.

\section{Effective Address Calculation}

For all memory operations, the effective address (EA) is calculated as:

\begin{equation}
	\text{EA} = \text{AddressRegister.high} : (\text{AddressRegister.low} + \text{Displacement})
\end{equation}

Where:
\begin{itemize}
	\item Address Register is one of: \reg{l}, \reg{a}, \reg{s}, or \reg{p}
	\item Displacement is either a 16-bit signed immediate value or the unshifted contents of the offset register
	      \texttt{rO}
\end{itemize}

The high word supplies the segment and is not changed by the displacement; only the low word participates in the
addition, modulo $2^{16}$.

\section{Legal Forms}

\begin{table}[H]
	\centering
	\small
	\begin{tabularx}{\textwidth}{lll>{\raggedright\arraybackslash}X}
		\toprule
		\textbf{Instruction} & \textbf{Syntax}                        & \textbf{Encoding} & \textbf{Notes}                                                  \\
		\midrule
		LOAD                 & \texttt{LOAD rD, (\#offset, addr)}     & \opcode{14}       & Immediate offset                                                \\
		LOAD                 & \texttt{LOAD rD, (rO, addr)}           & \opcode{15}       & Register offset; no shift suffix                                \\
		LOAD                 & \texttt{LOAD rD, (\#offset, addr)+}    & \opcode{16}       & Immediate offset; post-increment                                \\
		LOAD                 & \texttt{LOAD rD, (rO, addr)+}          & \opcode{17}       & Register offset; no shift suffix; post-increment                \\
		STOR                 & \texttt{STOR (\#offset, addr), rS}     & \opcode{10}       & Immediate offset                                                \\
		STOR                 & \texttt{STOR (rO, addr), rS[, shift]}  & \opcode{11}       & Register offset; optional source shift                          \\
		STOR                 & \texttt{STOR -(\#offset, addr), rS}    & \opcode{12}       & Immediate offset; pre-decrement                                 \\
		STOR                 & \texttt{STOR -(rO, addr), rS[, shift]} & \opcode{13}       & Register offset; optional shift; pre-decrement                  \\
		\bottomrule
	\end{tabularx}
	\caption{Legal Memory Instruction Forms}
	\label{tab:memory-legal-forms}
\end{table}

\noindent\textbf{Zero displacement:} Immediate-displacement memory forms may omit \texttt{\#0}. Thus
\texttt{LOAD rD, (addr)}, \texttt{LOAD rD, (addr)+}, \texttt{STOR (addr), rS}, and \texttt{STOR -(addr), rS} assemble
as the matching \texttt{\#0} forms. This shorthand does not add new addressing modes: \mnemonic{LOAD} has no
pre-decrement form and \mnemonic{STOR} has no post-increment form.

\noindent\textbf{Aliased register writes:} In post-increment \mnemonic{LOAD} forms the destination register must not
be either half of the auto-updated address-register pair; in pre-decrement \mnemonic{STOR} forms the source register
must not be either half of the auto-updated pair. The result of such an encoding is architecturally undefined and the
assembler rejects it.

\section{Common Memory Semantics}

\begin{table}[H]
	\centering
	\begin{tabular}{lp{9cm}}
		\toprule
		\textbf{Topic}      & \textbf{Definition}                                                                                                                                                                                                                                                                                                                                                                                                                                      \\
		\midrule
		Condition false     & The instruction is skipped and has no side effects. Registers, memory, address registers, and status flags are preserved.                                                                                                                                                                                                                                                                                                                                \\
		Address calculation & Immediate-offset forms add the 16-bit signed displacement to the selected address register low word. Register-offset forms add the unshifted \texttt{rO} to the selected address register low word; \texttt{rO} is never shifted, even on \mnemonic{STOR}'s register-offset forms, which shift the value being stored instead. The selected address register high word supplies the segment and is unchanged.                                                                                                                                                            \\
		Post-increment LOAD & The memory address is calculated from the original address register value. The address register low word is incremented by one word in write-back after the memory access.                                                                                                                                                                                                                                                                               \\
		Pre-decrement STOR  & The memory address uses the selected address register low word decremented by one word, plus the offset. The address register low word is decremented by one word in write-back after the memory access.                                                                                                                                                                                                                                                 \\
		Status flags        & Memory instructions preserve all status flags. Shifts used by memory instructions cannot update the status register.                                                                                                                                                                                                                                                                                                                                     \\
		Timing              & Memory instructions complete in 6 cycles when the bus access is acknowledged without wait states. Bus wait states extend the memory-access phase.                                                                                                                                                                                                                                                                                                        \\
		Exceptions          & Memory instructions can raise data bus, data bus-protection, segment-overflow, and privilege-violation faults; see the notes below.                                                                                                                                                                                                                                                                                                                      \\
		Fault side effects  & Destination register writes and address-register auto-updates occur in write-back after the memory-access phase, so a fault that aborts execution before write-back leaves them undone. Data bus, data bus-protection, segment-overflow, and privilege-violation faults are all retryable: the saved return address is the address of the faulting instruction, so returning from the fault handler with \mnemonic{RETE} retries the same memory access (see Chapter~\ref{ch:exceptions}). \\
		\bottomrule
	\end{tabular}
	\caption{Common Memory Instruction Semantics}
	\label{tab:memory-common-semantics}
\end{table}

\noindent Memory instructions can raise data bus and data bus-protection faults during the memory-access phase. The
16-bit displacement is signed two's-complement; effective-address calculation can raise a segment-overflow fault when
\texttt{SR.A} (Trap on Address Overflow) is set and the low-word address calculation wraps. This segment-overflow
fault is retryable: the saved return address is the address of the faulting instruction (see
Chapter~\ref{ch:exceptions}). A memory instruction raises a privilege-violation fault in
protected mode when its register field names a privileged register (\reg{sr}, \reg{ah}, \reg{lh}, \reg{ph}, or
\reg{sh}); this applies to the \mnemonic{LOAD} destination register and to the \mnemonic{STOR} source register, even
though \mnemonic{STOR} only reads it. Post-increment and pre-decrement address-register auto-update writes only the
low word of the pair, so it never changes the high word and never raises a privilege-violation fault. Documented
memory forms do not raise invalid-opcode faults. Instruction fetch can still raise the normal fetch faults before a
memory instruction executes. If trace mode was enabled when the instruction began, an instruction-trace fault is
raised after the instruction commits. Processing-unit encodings outside the documented memory forms are reserved or
architecturally undefined; they are not required to raise an invalid-opcode fault.

\section{Status Flag Effects}

Memory instructions do not update status flags. All four condition flags are preserved regardless of addressing form,
shift result, or auto-update side effect. This applies only to the memory-access \mnemonic{LOAD} forms
(\opcode{14}--\opcode{17}) and \mnemonic{STOR} (\opcode{10}--\opcode{13}); see Chapter~\ref{ch:alu-instructions} for
the flag behavior of the LOAD-immediate (\opcode{07}) and LOAD-register (\opcode{37}) forms. The AF field (bits 5--4)
is not a status-update source for these opcodes; it selects the address-register pair used to form the effective
address.

\begin{table}[H]
	\centering
	\begin{tabular}{lllllp{5.5cm}}
		\toprule
		\textbf{Instruction}            & \textbf{N} & \textbf{Z} & \textbf{C} & \textbf{V} & \textbf{Meaning}     \\
		\midrule
		LOAD (\opcode{14}--\opcode{17}) & -          & -          & -          & -          & Flags are preserved. \\
		STOR (\opcode{10}--\opcode{13}) & -          & -          & -          & -          & Flags are preserved. \\
		\bottomrule
	\end{tabular}
	\caption{Memory Instruction Status Flag Effects}
	\label{tab:memory-status-flag-effects}
\end{table}

See Table~\ref{tab:flag-effect-symbols} for the flag-effect symbol definitions.

\begin{instructionbox}{LOAD}{Load from Memory}
	\textbf{Applicability:} SIRC-1, all revisions

	\textbf{Opcodes:} \opcode{14} (Indirect Imm), \opcode{15} (Indirect Reg), \opcode{16} (Post-Inc Imm), \opcode{17} (Post-Inc Reg)

	\textbf{Instruction Fields:}
	\begin{itemize}
		\item \textbf{Register (bits 25--22, Immediate format, \opcode{14}/\opcode{16}):} Destination register
		      \texttt{rD}.
		\item \textbf{R1 (bits 25--22, Register format, \opcode{15}/\opcode{17}):} Destination register \texttt{rD}.
		\item \textbf{R2 (bits 21--18, Register format, \opcode{15}/\opcode{17}):} Unused, fixed at \texttt{0}. The
		      assembler does not accept a shift suffix on any \mnemonic{LOAD} form: \mnemonic{LOAD} has no field that
		      shifts the offset register or the loaded value. Use \mnemonic{SHFT} on \texttt{rO} before the
		      \mnemonic{LOAD} if a scaled offset is needed.
		\item \textbf{R3 (bits 17--14, Register format, \opcode{15}/\opcode{17}):} Offset register \texttt{rO}, used
		      unshifted.
		\item \textbf{Immediate (bits 21--6, Immediate format, \opcode{14}/\opcode{16}):} 16-bit signed displacement.
		\item \textbf{AF (bits 5--4):} Address-register pair -- \texttt{00} = \reg{l}, \texttt{01} = \reg{a},
		      \texttt{10} = \reg{s}, \texttt{11} = \reg{p}. Not a status-update source for these opcodes.
		\item \textbf{Condition field (bits 3--0):} All 16 condition-field encodings are legal; see \textbf{Condition
			      field} below.
	\end{itemize}

	\textbf{Syntax:}
	\begin{lstlisting}
LOAD rD, (#offset, addr)      ; rD = memory[addr + offset]
LOAD rD, (addr)               ; equivalent to LOAD rD, (#0, addr)
LOAD rD, (rO, addr)           ; rD = memory[addr + rO]
LOAD rD, (#offset, addr)+     ; rD = memory[addr + offset]; addr += 1
LOAD rD, (addr)+              ; equivalent to LOAD rD, (#0, addr)+
LOAD rD, (rO, addr)+          ; rD = memory[addr + rO]; addr += 1
\end{lstlisting}

	\textbf{Operands:} \texttt{rD} is the destination general-purpose register; a privileged register (\reg{sr},
	\reg{ah}, \reg{lh}, \reg{ph}, or \reg{sh}) raises a privilege-violation fault in protected mode. \texttt{addr} is
	the address-register pair that supplies the segment and base low word. Immediate-offset forms use a 16-bit signed
	displacement. Register-offset forms use \texttt{rO} as the displacement register, always unshifted; \mnemonic{LOAD}
	does not accept a shift suffix.

	\textbf{Operation:}
	\begin{verbatim}
if post_increment then
    effective_address = addr + offset
    destination = memory[effective_address]
    addr.low = addr.low + 1
else
    effective_address = addr + offset
    destination = memory[effective_address]
endif
; Flags are preserved
\end{verbatim}

	\textbf{Description:}

	Loads a 16-bit value from memory into the destination register. Post-increment forms update the address register after
	the load, useful for array traversal.

	\textbf{Flags:} Preserved.

	\textbf{Write-back:} If the condition is true, writes the loaded value to \texttt{rD}. Post-increment forms also write
	the incremented address-register low word.

	\textbf{Exceptions:} May raise data bus or data bus-protection faults during the memory read. May raise a
	segment-overflow fault during effective-address calculation when \texttt{SR.A} is set; this fault is not
	retryable (see Chapter~\ref{ch:exceptions}). May raise a privilege-violation fault in protected mode if
	\texttt{rD} names a privileged register.

	\textbf{Condition field:} If the condition is false, no memory read occurs, no destination register is written, no
	address register is incremented, and status flags are preserved.

	\textbf{Timing:} 6 cycles when the memory read is acknowledged without wait states; wait states extend the memory-access
	phase.

	\textbf{Privilege:} Available in protected mode and supervisor mode. Protected-mode data-memory access is subject to
	bus protection. A privilege-violation fault is raised in protected mode if \texttt{rD} names a privileged register
	(\reg{sr}, \reg{ah}, \reg{lh}, \reg{ph}, or \reg{sh}).

	\textbf{Example 14-1:}
	\begin{lstlisting}
; Load from fixed offset
LOAD r1, (#4, a)              ; r1 = memory[a + 4]

; Load using variable offset
LOAD r2, (r3, a)              ; r2 = memory[a + r3]

; Array iteration
LOAD r4, (#0, a)+             ; r4 = *a; a += 1

; Stack pop
LOAD r5, (#0, s)+             ; r5 = *s; s += 1

; LOAD has no shift suffix: scale the offset register first with SHFT
SHFT r2, LSL #2               ; r2 = r2 << 2
LOAD r3, (r2, a)               ; r3 = memory[a + r2]
\end{lstlisting}

\end{instructionbox}

\begin{instructionbox}{STOR}{Store to Memory}
	\textbf{Applicability:} SIRC-1, all revisions

	\textbf{Opcodes:} \opcode{10} (Indirect Imm), \opcode{11} (Indirect Reg), \opcode{12} (Pre-Dec Imm), \opcode{13} (Pre-Dec Reg)

	\textbf{Instruction Fields:}
	\begin{itemize}
		\item \textbf{Register (bits 25--22, Immediate format, \opcode{10}/\opcode{12}):} Source register \texttt{rS}.
		\item \textbf{R1 (bits 25--22, Register format, \opcode{11}/\opcode{13}):} Unused, encoded as 0.
		\item \textbf{R2 (bits 21--18, Register format, \opcode{11}/\opcode{13}):} Source register \texttt{rS}.
		\item \textbf{R3 (bits 17--14, Register format, \opcode{11}/\opcode{13}):} Offset register \texttt{rO}.
		\item \textbf{Immediate (bits 21--6, Immediate format, \opcode{10}/\opcode{12}):} 16-bit signed displacement.
		\item \textbf{Shift Type (bits 12--10, Register format, \opcode{11}/\opcode{13}):} \texttt{000} none,
		      \texttt{001} LSL, \texttt{010} LSR, \texttt{011} ASL, \texttt{100} ASR, \texttt{101} RTL, \texttt{110}
		      RTR, \texttt{111} reserved. Applies to \texttt{rS} before the value is written to memory.
		\item \textbf{Shift Amount (bits 9--6, Register format, \opcode{11}/\opcode{13}):} Immediate count 0--15, or
		      (Shift Operand bit set) a register holding the count; a register count above 15 is architecturally
		      undefined.
		\item \textbf{AF (bits 5--4):} Address-register pair -- \texttt{00} = \reg{l}, \texttt{01} = \reg{a},
		      \texttt{10} = \reg{s}, \texttt{11} = \reg{p}. Not a status-update source for these opcodes.
		\item \textbf{Condition field (bits 3--0):} All 16 condition-field encodings are legal; see \textbf{Condition
			      field} below.
	\end{itemize}

	\textbf{Syntax:}
	\begin{lstlisting}
STOR (#offset, addr), rS      ; memory[addr + offset] = rS
STOR (addr), rS               ; equivalent to STOR (#0, addr), rS
STOR (rO, addr), rS           ; memory[addr + rO] = rS
STOR -(#offset, addr), rS     ; addr -= 1; memory[addr + offset] = rS
STOR -(addr), rS              ; equivalent to STOR -(#0, addr), rS
STOR -(rO, addr), rS          ; addr -= 1; memory[addr + rO] = rS

; With optional shift on the source value (register-offset memory forms only)
STOR (rO, addr), rS, shift           ; Shift rS, then store
STOR -(rO, addr), rS, shift          ; Decrement, shift rS, then store
\end{lstlisting}

	\textbf{Operands:} \texttt{rS} is the source general-purpose register; a privileged register (\reg{sr}, \reg{ah},
	\reg{lh}, \reg{ph}, or \reg{sh}) raises a privilege-violation fault in protected mode, even though \mnemonic{STOR}
	only reads it. \texttt{addr} is the address-register pair that supplies the segment and base low word.
	Immediate-offset forms use a 16-bit signed displacement. Register-offset forms use \texttt{rO} as the displacement
	register and may apply a shift to the source value before storing it.

	\textbf{Operation:}
	\begin{verbatim}
if pre_decrement then
    decremented_low = addr.low - 1
    effective_address = (addr.high : decremented_low) + offset
    memory[effective_address] = shift_if_encoded(source)
    addr.low = decremented_low
else
    effective_address = addr + offset
    memory[effective_address] = shift_if_encoded(source)
endif
; Source register and flags are preserved
\end{verbatim}

	\textbf{Description:}

	Stores a 16-bit value from a register into memory. Pre-decrement forms use the decremented address for the store and
	write the decremented address register value back after the memory access, useful for stack push operations.

	\textbf{Flags:} Preserved.

	\textbf{Write-back:} If the condition is true, writes the source value to memory. Pre-decrement forms also write the
	decremented address-register low word.

	\textbf{Exceptions:} May raise data bus or data bus-protection faults during the memory write. May raise a
	segment-overflow fault during effective-address calculation when \texttt{SR.A} is set; this fault is not
	retryable (see Chapter~\ref{ch:exceptions}). May raise a privilege-violation fault in protected mode if
	\texttt{rS} names a privileged register.

	\textbf{Condition field:} If the condition is false, no memory write occurs, no address register is decremented, and
	status flags are preserved.

	\textbf{Timing:} 6 cycles when the memory write is acknowledged without wait states; wait states extend the memory-access
	phase.

	\textbf{Privilege:} Available in protected mode and supervisor mode. Protected-mode data-memory access is subject to
	bus protection. A privilege-violation fault is raised in protected mode if \texttt{rS} names a privileged register
	(\reg{sr}, \reg{ah}, \reg{lh}, \reg{ph}, or \reg{sh}), even though \mnemonic{STOR} only reads it.

	\textbf{Example 14-2:}
	\begin{lstlisting}
; Store to fixed offset
STOR (#8, a), r1              ; memory[a + 8] = r1

; Store using variable offset
STOR (r3, a), r2              ; memory[a + r3] = r2

; Store with shift (data shifted before storing)
STOR (r2, a), r1, LSL #1      ; memory[a + r2] = r1 << 1
STOR (r3, a), r1, LSR #3      ; memory[a + r3] = r1 >> 3

; Stack push
STOR -(#0, s), r4             ; s -= 1; memory[s] = r4
\end{lstlisting}

\end{instructionbox}

\section{Common Memory Access Patterns}

\subsection{Structure Field Access}
\begin{lstlisting}
; struct { int16_t x, y, z; } point;
; 'a' points to point
LOAD r1, (#0, a)              ; r1 = point.x
LOAD r2, (#1, a)              ; r2 = point.y
LOAD r3, (#2, a)              ; r3 = point.z
\end{lstlisting}

\subsection{Array Iteration}
\begin{lstlisting}
; Process array of 16-bit values
; 'a' = array base, r7 = count
:loop
    LOAD r1, (#0, a)+         ; Load element, advance
    ; ... process r1 ...
    SUBI r7, #1
    BRAN|!= @loop
\end{lstlisting}

\subsection{Stack Operations}
\begin{lstlisting}
; Push registers
STOR -(#0, s), r1
STOR -(#0, s), r2
STOR -(#0, s), r3

; Pop registers (reverse order)
LOAD r3, (#0, s)+
LOAD r2, (#0, s)+
LOAD r1, (#0, s)+
\end{lstlisting}

\section{Shift Operations with Memory Instructions}
\label{sec:memory-shifts}

Only \mnemonic{STOR}'s register-offset forms support an optional shift suffix, applied to the source register value
before it is written to memory. \mnemonic{LOAD} does not support a shift suffix on any form -- the offset register
\texttt{rO} is always used unshifted, and there is no field that shifts the loaded value. This asymmetry is
architectural, not an oversight: use \mnemonic{SHFT} on the offset register before a \mnemonic{LOAD} if a scaled
index is needed.

\subsection{STOR with Shift}

When storing to memory, the source register value is shifted \textbf{before} being written to memory. The original
register value is \textbf{not} modified:

\begin{lstlisting}
STOR (r2, a), r1, LSL #1      ; memory[a + r2] = r1 << 1 (r1 unchanged)
STOR (r2, a), r1, ASL #2      ; memory[a + r2] = r1 << 2 (r1 unchanged)
STOR (r4, a), r1, ASR #4      ; memory[a + r4] = r1 >> 4 (r1 unchanged)
\end{lstlisting}

\subsection{Supported Shift Types}

All standard shift types are supported:
\begin{itemize}
	\item \textbf{LSL} -- Logical Shift Left (fill with zeros)
	\item \textbf{LSR} -- Logical Shift Right (fill with zeros)
	\item \textbf{ASL} -- Arithmetic Shift Left (same data result as LSL; sets V on signed overflow)
	\item \textbf{ASR} -- Arithmetic Shift Right (sign-extend)
	\item \textbf{RTL} -- Rotate Left within the operand
	\item \textbf{RTR} -- Rotate Right within the operand
\end{itemize}

\subsection{Shift Limitations}

\textbf{Important:} The shift suffix is \textbf{only} supported on \mnemonic{STOR}'s register-offset forms. The
following do \textbf{not} support a shift suffix:
\begin{itemize}
	\item Any \mnemonic{LOAD} form, register-offset included -- \texttt{LOAD r1, (r2, a)} -- No shift suffix exists
	\item Loading immediate values: \texttt{LOAD r1, \#0x1234} -- No shift allowed
	\item Copying registers: \texttt{LOAD r1, r2} -- No shift allowed
	\item Immediate-offset memory load: \texttt{LOAD r1, (\#0, a)} -- No shift allowed
	\item Immediate-offset memory store: \texttt{STOR (\#0, a), r1} -- No shift allowed
\end{itemize}

For any of these, load or store the unshifted value and use an ALU instruction or \mnemonic{SHFT} for the shift
operation.

\subsection{Use Cases}

\subsubsection{Scaled Indexing}

\mnemonic{LOAD} has no shift suffix, so scaling an index for array access takes a separate \mnemonic{SHFT} before the
\mnemonic{LOAD}:

\begin{lstlisting}
; r2 holds an element index; each element occupies 4 words
SHFT r2, LSL #2                ; r2 = r2 << 2
LOAD r3, (r2, a)               ; r3 = memory[a + r2]
\end{lstlisting}

\subsubsection{Bit Manipulation}

Shifting the source register before a \mnemonic{STOR} is useful for packing a scaled or shifted value into memory
without a separate ALU instruction:

\begin{lstlisting}
; Store a scaled value
STOR (r3, a), r2, LSL #2      ; memory[a + r3] = r2 * 4
\end{lstlisting}

\subsubsection{Combined with Auto-Increment/Decrement}

\mnemonic{STOR}'s shift suffix works seamlessly with pre-decrement mode:

\begin{lstlisting}
; Decrement, shift the source value, and store
STOR -(r2, a), r1, LSR #2     ; a -= 1; memory[a + r2] = r1 >> 2
\end{lstlisting}

\subsection{Status Flags}

\textbf{Note:} Unlike shifts in ALU instructions, a shift performed during a \mnemonic{STOR} cannot update the status
register flags. For flag updates based on a shift result, use an ALU instruction or the \mnemonic{SHFT}
meta-instruction (see Chapter~\ref{ch:alu-instructions}).
